% Date : 26/12/2022
% Version : 
% Auteur : CRG
% Relu :  19/02/2022
% Relecteur : Lucas Henry
% Harmonisateur : Jimmy Roussel


% ------------------------------------------------------------ %
% ----------------------  Fiche CHI03  ----------------------- %
% Thème : Chimie
% Classes : Toutes classes de 1ère année
% ------------------------------------------------------------ %



% ======================================================= % 
% ============== Gestion de la compilation ============== %
% ======================================================= % 
\ifdefined\mainIsLoaded\else
\RequirePackage{import}
\subimport{../../}{_preambule_CdE_PC}
\begin{document}
\fi
% ======================================================= % 
% ======================================================= % 



% ======================================================= % 
% ===================== Couleurs ======================== %
% ======================================================= % 

% -----------------------------------%
% Merci de préfixer vos couleur

% La commande \couleurNBouCouleur prend deux paramètres :
% #1 est la couleur NB ; #2 est la couleur "en couleur"

\def\CBDcouleurA{\couleurNBouCouleur{black}{blue}}

% -----------------------------------%

% ======================================================= % 
% ======================================================= % 




% ============================================================================ %
%                            FICHE D'ENTRAÎNEMENT                              %
% ============================================================================ %
% ======================= Métadonnées sur la fiche =========================== %
\titreFicheEntrainement		{Réactions chimiques}
\grandTheme					{Chimie}
\numeroFiche				{CHI03}
\uniqueID					{FvsGjyRGSH} % une chaîne de caractères (a-z, A-Z) aléatoire de 10 caractères.
\nombreColonnesReponses		{1} % 1, 2, 3, etc.
% ============================================================================ %

%%%%%%%%%%%%%%%%%%%%%%%%%%%%%%%%%%%%%%%%%%%%%%%%%%%%%%%%%%%%%%%%%%%%%%%%%%%%%%%%%%%%%%%%%%%%%%
\debutFicheEntrainement                                                                      %
%%%%%%%%%%%%%%%%%%%%%%%%%%%%%%%%%%%%%%%%%%%%%%%%%%%%%%%%%%%%%%%%%%%%%%%%%%%%%%%%%%%%%%%%%%%%%%




% --------------------------------------------------------------------- %
%                              Prérequis                                %
%                             (facultatif)                              %
% --------------------------------------------------------------------- %
% ================== Métadonnées sur le prérequis ===================== %
\avecPrerequis{Y} % Y ou N
% ===================================================================== %
\begin{prerequis}
Tableaux d'avancement, avancement ($\xi$) et avancement volumique ($\xi_v$) d'une réaction. Loi d'action de masse. Définition du pH, constante d'acidité. Constante d'autoprotolyse de l'eau.
\end{prerequis}
%--------------------------------------------------------------------- %






% •••••••••••••••••••••••••••••••••••••••••••••••••••••••••••••••••••••••••••• %
% •••••••••••••••••••••••••••••••••••••••••••••••••••••••••••••••••••••••••••• %
\sectionFicheEntrainement{Pour commencer}
% •••••••••••••••••••••••••••••••••••••••••••••••••••••••••••••••••••••••••••• %
% •••••••••••••••••••••••••••••••••••••••••••••••••••••••••••••••••••••••••••• %







% ********************************************************************* %
%                           ENTRAÎNEMENT      1                          % 
% ********************************************************************* %
% ================ Métadonnées sur l'entraînement ===================== %

\titreEntrainementFacultatif	{Ajuster des équations de réaction}
\hauteurLargeurCadreReponse		{8mm}{7.5cm}
\dureeResolutionFacultative		{2} % 1, 2, 3 ou 4
\typeEntrainement				{Newton} % Newton ou QCM ou AN ou vide (exercice "normal")
\basiqueEtTransversal			{Y}
\nombreColonnesQuestions		{1} % vide, 1, 2, 3, etc.
\avecPlusieursQuestions			{Y} % Y ou N
\initialisationEntrainement
% ===================================================================== %

Ajuster les équations des réactions suivantes.

% =============================================================================== %
\debutEntrainement
% =============================================================================== %

% ^^^^^^^^^^^^^^^^^^^^^^^^^^^^^^^^^^^^^^ %
%               Question    a             %
% ^^^^^^^^^^^^^^^^^^^^^^^^^^^^^^^^^^^^^^ %

\begin{enonce}
CO + O$_2$ = CO$_2$
\end{enonce}

\reponse{$2\,\mathrm{CO}+\mathrm{O_2}=2\,\mathrm{CO_2}$}

\begin{corrige}
On commence d'abord par équilibrer les atomes de carbone (un de chaque côté). On a deux atomes d'oxygène à droite, on doit donc en placer deux à gauche. Ce qui donne : $\mathrm{CO}+\frac12\,\mathrm{O_2}=\mathrm{CO_2}$.
 
On préfère raisonner avec des coefficients st\oe{}chiométriques entiers, il suffit alors de multiplier les coefficients par deux : $2\,\mathrm{CO}+\mathrm{O_2}=2\,\mathrm{CO_2}$. 

\end{corrige}

% ************************************** %


% ^^^^^^^^^^^^^^^^^^^^^^^^^^^^^^^^^^^^^^ %
%               Question   b              %
% ^^^^^^^^^^^^^^^^^^^^^^^^^^^^^^^^^^^^^^ %

\begin{enonce}
Ag$^+$ + Cu = Ag + Cu$^{2+}$
\end{enonce}

\reponse{2 Ag$^+$ + Cu = 2 Ag + Cu$^{2+}$}

\begin{corrige}
Initialement, les charges ne sont pas équilibrées. Il faut mettre 2 Ag$^+$ pour ajuster les charges. Enfin, on équilibre l'élément Ag en mettant un coefficient 2 au produit Ag. On obtient 2 Ag$^+$ + Cu = 2 Ag + Cu$^{2+}$.
\end{corrige}

% ************************************** %


% ^^^^^^^^^^^^^^^^^^^^^^^^^^^^^^^^^^^^^^ %
%               Question     c            %
% ^^^^^^^^^^^^^^^^^^^^^^^^^^^^^^^^^^^^^^ %

\begin{enonce}
NO + CO = N$_2$ + CO$_2$
\end{enonce}

\reponse{2 NO + 2 CO = N$_2$ + 2 CO$_2$}

\begin{corrige}
On commence par équilibrer l'élément azote : 2~NO~+~CO~=~N$_2$~+~CO$_2$. 
Les carbones sont équilibrés mais pas les atomes d'oxygène. On doit donc trouver $x$ tel que : \minusMedskip
\[
2\,\mathrm{NO}+x\,\mathrm{CO}=\mathrm{N_2}+x\,\mathrm{CO_2}. \minusMedskip
\]
En raisonnant sur l'atome d'oxygène on trouve $2+x=2x$ soit $x=2$.
\end{corrige}

% ************************************** %



% ^^^^^^^^^^^^^^^^^^^^^^^^^^^^^^^^^^^^^^ %
%               Question     d            %
% ^^^^^^^^^^^^^^^^^^^^^^^^^^^^^^^^^^^^^^ %

\begin{enonce}
S$_2$O$^{2-}_8$ + I$^-$ = SO$^{2-}_4$ + I$_2$
\end{enonce}

\reponse{$\mathrm{S_2O^{2-}_8}+2\,\mathrm{I^-}=2\,\mathrm{SO^{2-}_4}+\mathrm{I_2}$}

\begin{corrige}
Commençons par équilibrer les atomes d'iode puis le soufre et enfin l'oxygène. On arrive à : \minusMedskip
\[
\mathrm{S_2O^{2-}_8}+2\,\mathrm{I^-}=2\,\mathrm{SO^{2-}_4}+\mathrm{I_2}. \minusMedskip
\]
On s'aperçoit que les charges sont \emph{de facto} ajustées. La réaction est équilibrée !
\end{corrige}

% ************************************** %


% ^^^^^^^^^^^^^^^^^^^^^^^^^^^^^^^^^^^^^^ %
%               Question    e             %
% ^^^^^^^^^^^^^^^^^^^^^^^^^^^^^^^^^^^^^^ %

\begin{enonce}
C$_8$H$_{18}$ + O$_2$ = CO$_2$ + H$_2$O
\end{enonce}

\reponse{$\mathrm{2\,C_8H_{18}}+25\,\mathrm{O_2}=16\,\mathrm{CO_2}+18\,\mathrm{H_2O}$}

\begin{corrige}
Commençons par ajuster les atomes d'hydrogène : $\mathrm{C_8H_{18}}+\mathrm{O_2}=\mathrm{CO_2}+9\,\mathrm{H_2O}$.
Poursuivons avec les atomes de  carbone : $\mathrm{C_8H_{18}}+\mathrm{O_2}=8\mathrm{CO_2}+9\,\mathrm{H_2O}$.
Puis avec les atomes d'oxygène : $\mathrm{C_8H_{18}}+\frac{25}{2}\mathrm{O_2}=8\mathrm{CO_2}+9\,\mathrm{H_2O}$.

Terminons en multipliant tous les coefficients par deux : $\mathrm{2\,C_8H_{18}}+25\,\mathrm{O_2}=16\,\mathrm{CO_2}+18\,\mathrm{H_2O}$.
\end{corrige}

% ************************************** %


% ^^^^^^^^^^^^^^^^^^^^^^^^^^^^^^^^^^^^^^ %
%               Question      f           %
% ^^^^^^^^^^^^^^^^^^^^^^^^^^^^^^^^^^^^^^ %

\begin{enonce}
MnO$^-_4$ + H$^+$ + Fe$^{2+}$ = Fe$^{3+}$ + Mn$^{2+}$ + H$_2$O
\end{enonce}

\reponse{$\mathrm{MnO^-_4}+8\,\mathrm{H^+}+5\,\mathrm{Fe^{2+}}=5\,\mathrm{Fe^{3+}}+\mathrm{Mn^{2+}}+4\, \mathrm{H_2O}$}

\begin{corrige}
Commençons par équilibrer les atomes d'oxygène : 
$\mathrm{MnO^-_4}+\mathrm{H^+}+\mathrm{Fe^{2+}}=\mathrm{Fe^{3+}}+\mathrm{Mn^{2+}}+4\, \mathrm{H_2O}$.

Puis les atomes d'hydrogène : $\mathrm{MnO^-_4}+8\,\mathrm{H^+}+\mathrm{Fe^{2+}}=\mathrm{Fe^{3+}}+\mathrm{Mn^{2+}}+4\, \mathrm{H_2O}$.

Les éléments sont équilibrés. Comptons les charges : $+9$ à gauche et $+5$ à droite. Les charges ne sont donc pas ajustées. Or, on n'a pas encore considéré le fer. Appelons $x$ son coefficient :  
\[
\mathrm{MnO^-_4}+8\,\mathrm{H^+}+x\,\mathrm{Fe^{2+}}=x\,\mathrm{Fe^{3+}}+\mathrm{Mn^{2+}}+4\, \mathrm{H_2O}.
\]
L'équilibre des charges  donne $7+2x=2+3x$, d'où $x=5$.
\end{corrige}

% ************************************** %


% =============================================================================== %
\finEntrainement
% =============================================================================== %





% ********************************************************************* %
%                           ENTRAÎNEMENT   2                             % 
% ********************************************************************* %
% ================ Métadonnées sur l'entraînement ===================== %

\titreEntrainementFacultatif	{Tableau d'avancement}
\hauteurLargeurCadreReponse		{7mm}{35mm}
\dureeResolutionFacultative		{1} % 1, 2, 3 ou 4
\typeEntrainement				{Newton} % Newton ou QCM ou AN ou vide (exercice "normal")
\nombreColonnesQuestions		{1} % vide, 1, 2, 3, etc.
\avecPlusieursQuestions			{N} % Y ou N
\initialisationEntrainement
% ===================================================================== %

On considère le tableau d'avancement en quantité de matière suivant :
\begin{center}
\renewcommand{\arraystretch}{1.5}
\begin{tabular}{|l|c|c||c|}
\hline
&
%
    \multicolumn{1}{c@{\protect\makebox[0cm]{$+$}}}{%
      ~~~~$\mathrm{{N_2}_\text{(g)}}$~~~~}                   &
    \multicolumn{1}{c@{\protect\makebox[0cm]{$=$}}}{%
      ~~~~$\mathrm{3~{H_2}_\text{(g)}}$~~~~}                        &
        \multicolumn{1}{c|}{%
      ~~~~$\mathrm{2~{NH_3}_\text{(g)}}$~~~}              \\ \hline
%
  \textbf{\'Etat initial}        & $n_1$                       &
    $n_2$                 & 0                                                    \\ \hline
%
  \textbf{\'Etat final}	&$\alpha$	&$\beta$	&$\gamma$	\\ \hline
\end{tabular}
\end{center}

où $n_1$ et $n_2$ sont des quantités de matière. À l'instant final, l'avancement molaire de la réaction vaut $\xi$. 

\medskip

Déterminer en fonction de  $n_1$, $n_2$ et $\xi$, les quantités suivantes : 


% =============================================================================== %
\debutEntrainement
% =============================================================================== %

% ^^^^^^^^^^^^^^^^^^^^^^^^^^^^^^^^^^^^^^ %
%               Question                 %
% ^^^^^^^^^^^^^^^^^^^^^^^^^^^^^^^^^^^^^^ %

\begin{enonce}
$\alpha$
\end{enonce}

\reponse{$n_1-\xi$}

\begin{corrige}
Par définition l'avancement est lié aux quantités de matière des produits ou réactifs \emph{via} $\xi=\frac{n_i(t)-n_i(0)}{\nu_i}$ où $\nu_i$ est le coefficient st\oe chiométrique algébrique du produit ou réactif. On obtient donc : \minusMedskip
\begin{center}
\renewcommand{\arraystretch}{1.2}
\begin{tabular}{|l|c|c||c|}
  \hline
&
%
    \multicolumn{1}{c@{\protect\makebox[0cm]{$\mathrm{+}$}}}{%
      ~~~~$\mathrm{{N_2}_\text{(g)}}$~~~~}                   &
    \multicolumn{1}{c@{\protect\makebox[0cm]{$=$}}}{%
      ~~~~$\mathrm{3 {H_2}_\text{(g)}}$~~~~}                        &
        \multicolumn{1}{c|}{%
      ~~~~$\mathrm{2 {NH_3}_\text{(g)}}$~~~}              \\ \hline
%
  \textbf{\'Etat initial}        & $n_1$                       &
    $n_2$                 & 0                                                    \\ \hline
%
  \textbf{\'Etat final} & $n_1-\xi$                       &
    $n_2-3\xi$                 & $2\xi$                                          \\ \hline
\end{tabular}
\end{center}
\end{corrige}

% ************************************** %


% ^^^^^^^^^^^^^^^^^^^^^^^^^^^^^^^^^^^^^^ %
%               Question                 %
% ^^^^^^^^^^^^^^^^^^^^^^^^^^^^^^^^^^^^^^ %

\begin{enonce}
$\beta$
\end{enonce}

\reponse{$n_2-3\xi$}

% ************************************** %


% ^^^^^^^^^^^^^^^^^^^^^^^^^^^^^^^^^^^^^^ %
%               Question                 %
% ^^^^^^^^^^^^^^^^^^^^^^^^^^^^^^^^^^^^^^ %

\begin{enonce}
$\gamma$
\end{enonce}

\reponse{$2 \xi$}

% ************************************** %


% =============================================================================== %
\finEntrainement
% =============================================================================== %





\newpage



% ********************************************************************* %
%                           ENTRAÎNEMENT          3                      % 
% ********************************************************************* %
% ================ Métadonnées sur l'entraînement ===================== %
\titreEntrainementFacultatif	{Dimension de la constante thermodynamique d'équilibre}
\hauteurLargeurCadreReponse		{6mm}{1.5cm}
\dureeResolutionFacultative		{1} % 1, 2, 3 ou 4
\typeEntrainement				{QCM} % Newton ou QCM ou AN ou vide (exercice "normal")
\nombreColonnesQuestions		{1} % vide, 1, 2, 3, etc.
\avecPlusieursQuestions			{N} % Y ou N
\initialisationEntrainement
% ===================================================================== %


% =============================================================================== %
\debutEntrainement
% =============================================================================== %


% ^^^^^^^^^^^^^^^^^^^^^^^^^^^^^^^^^^^^^^ %
%               Question                 %
% ^^^^^^^^^^^^^^^^^^^^^^^^^^^^^^^^^^^^^^ %
\begin{enonce}
	On considère la transformation d'équation : 
	\begin{center}
	{SO$_2$Cl$_2$}$_\text{(g)}$ = {SO$_2$}$_\text{(g)}$ + {Cl$_2$}$_\text{(g)}$.
	\end{center}
	Trouver, parmi les formules suivantes, l'expression de sa constante d'équilibre $K^\circ$ :
	
	\medskip
	
	\begin{listeQCM2Colonnes}
	\item $K^\circ=\frac{P(\text{SO}_2)_\text{eq}\times P(\text{Cl}_2)_\text{eq}}{P(\text{SO}_2\text{Cl}_2)_\text{eq}}$\bigskip
	\item $K^\circ=\frac{P(\text{SO}_2\text{Cl}_2)_\text{eq}}{P(\text{SO}_2)_\text{eq}\times P(\text{Cl}_2)_\text{eq}}$
	\columnbreak
	\item $K^\circ=\frac{P(\text{SO}_2\text{Cl}_2)_\text{eq}\times P^\circ}{P(\text{SO}_2)_\text{eq}\times P(\text{Cl}_2)_\text{eq}}$\bigskip
	\item $K^\circ=\frac{P(\text{SO}_2)_\text{eq}\times P(\text{Cl}_2)_\text{eq}}{P(\text{SO}_2\text{Cl}_2)_\text{eq}\times P^\circ}$
	\end{listeQCM2Colonnes}
\smallskip

\end{enonce}

\reponse{\reponseD{}}

\begin{corrige}
La constante thermodynamique d'équilibre est une grandeur adimensionnée, ce qui exclut les propositions \reponseA{} et \reponseB{}. Ensuite, par définition, l'activité des produits de la réaction doit se trouver au numérateur et celle des réactifs au dénominateur. On garde donc l'expression \reponseD{}.
\end{corrige}

% ************************************** %

% =============================================================================== %
\finEntrainement
% =============================================================================== %







% ********************************************************************* %
%                           ENTRAÎNEMENT             4                   % 
% ********************************************************************* %
% ================ Métadonnées sur l'entraînement ===================== %
\titreEntrainementFacultatif	{Expression de la constante thermodynamique d'équilibre}
\hauteurLargeurCadreReponse		{6mm}{1.5cm}
\dureeResolutionFacultative		{1} % 1, 2, 3 ou 4
\typeEntrainement				{QCM} % Newton ou QCM ou AN ou vide (exercice "normal")
\nombreColonnesQuestions		{1} % vide, 1, 2, 3, etc.
\avecPlusieursQuestions			{N} % Y ou N
\initialisationEntrainement
% ===================================================================== %


% =============================================================================== %
\debutEntrainement
% =============================================================================== %


% ^^^^^^^^^^^^^^^^^^^^^^^^^^^^^^^^^^^^^^ %
%               Question                 %
% ^^^^^^^^^^^^^^^^^^^^^^^^^^^^^^^^^^^^^^ %
\begin{enonce}
	On considère la transformation d'équation : 
	\begin{center}
	{Cd(OH)$_2$}$_\text{(s)}$ + 4 {NH$_3$}$_\text{(aq)}$ = {[Cd(NH$_3$)$_4$]$^{2+}$}$_\text{(aq)}$ + 2 {HO$^-$}$_\text{(aq)}$.
	\end{center}
	Trouver, parmi les formules suivantes, l'expression de sa constante d'équilibre $K^\circ$ :
	
	\medskip
	
	\begin{listeQCM2Colonnes}
	\item $K^\circ=\frac{\left[\text{HO}^-\right]_\text{eq}\times\left[[\text{Cd(NH}_3)_4]^{2+}\right]_\text{eq}}{\left[\text{Cd(OH)}_2\right]_\text{eq}\times\left[\text{NH}_3\right]_\text{eq}}$\bigskip
	\item $K^\circ=\frac{\left[\text{HO}^-\right]^2_\text{eq}\times\left[[\text{Cd(NH}_3)_4]^{2+}\right]_\text{eq}}{\left[\text{Cd(OH)}_2\right]_\text{eq}\times\left[\text{NH}_3\right]^4_\text{eq}\times (C^\circ)^2}$\bigskip
	\item $K^\circ=\frac{\left[\text{HO}^-\right]^2_\text{eq}\times\left[[\text{Cd(NH}_3)_4]^{2+}\right]_\text{eq}\times (C^\circ)^2}{\left[\text{Cd(OH)}_2\right]_\text{eq}\times\left[\text{NH}_3\right]^4_\text{eq}}$\columnbreak
	\item $K^\circ=\frac{\left[\text{HO}^-\right]^2_\text{eq}\times\left[[\text{Cd(NH}_3)_4]^{2+}\right]_\text{eq}}{\left[\text{NH}_3\right]^4_\text{eq}\times C^\circ}$\bigskip
	\item $K^\circ=\frac{\left[\text{HO}^-\right]^2_\text{eq}\times\left[[\text{Cd(NH}_3)_4]^{2+}\right]_\text{eq}\times C^\circ}{\left[\text{NH}_3\right]^4_\text{eq}}$\bigskip
	\item $K^\circ=\frac{\left[\text{NH}_3\right]^4_\text{eq}\times C^\circ}{\left[\text{HO}^-\right]^2_\text{eq}\times\left[[\text{Cd(NH}_3)_4]^{2+}\right]_\text{eq}}$
	\end{listeQCM2Colonnes}

\bigskip

\end{enonce}

\reponse{\reponseE{}}

\begin{corrigeNewpage}
La constante thermodynamique d'équilibre est une grandeur adimensionnée, ce qui exclut les propositions \reponseB{}, \reponseD{} et \reponseF{}. Ensuite, par définition, l'activité d'un solide seul dans sa phase vaut 1, ce qui exclut les propositions \reponseA{} et \reponseC{}. On garde donc l'expression \reponseE{}.
\end{corrigeNewpage}

% ************************************** %

% =============================================================================== %
\finEntrainement
% =============================================================================== %







% ********************************************************************* %
%                           ENTRAÎNEMENT             5                   % 
% ********************************************************************* %
% ================ Métadonnées sur l'entraînement ===================== %

\titreEntrainementFacultatif	{Expression et calcul de la constante d'équilibre}
\hauteurLargeurCadreReponse		{11mm}{5cm}
\dureeResolutionFacultative		{2} % 1, 2, 3 ou 4
\typeEntrainement				{} % Newton ou QCM ou AN ou vide (exercice "normal")
\nombreColonnesQuestions		{1} % vide, 1, 2, 3, etc.
\avecPlusieursQuestions			{Y} % Y ou N
\initialisationEntrainement
% ===================================================================== %

On considère la réaction acide-base entre le chlorure d'ammonium ($\mathrm{NH^+_4}$ ; $\mathrm{Cl^-}$) et l'hydroxyde de sodium (Na$^+$ ; HO$^-$) :
\[
\mathrm{NH^+_4}_\text{(aq)}+\mathrm{HO^-}_\text{(aq)}=
\mathrm{NH_3}_\text{(aq)}+\mathrm{H_2O}_\text{($\ell$)}.
\]
% =============================================================================== %
\debutEntrainement
% =============================================================================== %

% ^^^^^^^^^^^^^^^^^^^^^^^^^^^^^^^^^^^^^^ %
%               Question      a           %
% ^^^^^^^^^^^^^^^^^^^^^^^^^^^^^^^^^^^^^^ %

\begin{enonce}
En utilisant la loi d'action de masse, exprimer la constante d'équilibre $K^\circ$ de la réaction en fonction des activités des différentes espèces physico-chimiques intervenant dans la réaction.

\end{enonce}


\reponse{$\frac{a(\text{NH}_3)_\text{eq}\times a(\text{H}_2\text{O})_\text{eq}}{a(\text{NH}^+_4)_\text{eq}\times a(\text{HO}^-)_\text{eq}}$}

\begin{corrige}
D'après la loi d'action de masse, on a  $K^\circ=Q_\text{eq}=\frac{a(\text{NH}_3)_\text{eq}\times a(\text{H}_2\text{O})_\text{eq}}{a(\text{NH}^+_4)_\text{eq}\times a(\text{HO}^-)_\text{eq}}$.
\end{corrige}

% ************************************** %



% ^^^^^^^^^^^^^^^^^^^^^^^^^^^^^^^^^^^^^^ %
%               Question     b           %
% ^^^^^^^^^^^^^^^^^^^^^^^^^^^^^^^^^^^^^^ %

\begin{enonce}
La constante d'acidité $K_A$ du couple NH$^+_4$/NH$_3$ est la constante d'équilibre de la réaction 
\[
\mathrm{NH_4^+}_\text{(aq)}+\mathrm{H_2O=NH_3}_\text{(aq)}+\mathrm{H_3O^+}_\text{(aq)}.\minusMedskip
\]
Exprimer $K_A$ en fonction des activités des espèces pertinentes
\end{enonce}

\reponse{$\frac{a(\text{NH}_3)_\text{eq}\times a(\text{H}_3\text{O}^+)_\text{eq}}{a(\text{NH}^+_4)_\text{eq}\times a(\text{H}_2\text{O})_\text{eq}}$}

\begin{corrige}
La constante d'acidité est la constante d'équilibre associée à la réaction entre l'acide du couple et l'eau : \minusMedskip
\[
\text{NH}^+_4 + \text{H}_2\text{O}=\text{NH}_3 + \text{H}_3\text{O}^+.
\]
D'après la loi d'action de masse, on a donc : $K_A=\frac{a(\text{NH}_3)_\text{eq}\times a(\text{H}_3\text{O}^+)_\text{eq}}{a(\text{NH}^+_4)_\text{eq}\times a(\text{H}_2\text{O})_\text{eq}}$.
\end{corrige}

% ************************************** %

% ^^^^^^^^^^^^^^^^^^^^^^^^^^^^^^^^^^^^^^ %
%               Question      c           %
% ^^^^^^^^^^^^^^^^^^^^^^^^^^^^^^^^^^^^^^ %

\begin{enonce}
La constante d'autoprotolyse de l'eau $K_e$ est la constante d'équilibre de la réaction
\[
2\,{\mathrm{H_2O}}_\text{($\ell$)}=\mathrm{H_3O^+}_\text{(aq)}+\mathrm{HO^-}_\text{(aq)}.\vspace{-3mm}
\]
Exprimer $K_e$ en fonction des activités des espèces pertinentes
\end{enonce}

\reponse{$\frac{a(\text{HO}^-)_\text{eq}\times a(\text{H}_3\text{O}^+)_\text{eq}}{a(\text{H}_2\text{O})^2_\text{eq}}$}

\begin{corrige}
La constante d'autoprotolyse de l'eau est la constante d'équilibre associée à la réaction : 
\[
2~\text{H}_2\text{O}=\text{H}_3\text{O}^+ + \text{HO}^-.
\]
D'après la loi d'action de masse, on a donc 
$K_e=\frac{a(\text{H}_3\text{O}^+)_\text{eq}\times a(\text{HO}^-)_\text{eq}}{a(\text{H}_2\text{O})^2_\text{eq}}$.
\end{corrige}

% ************************************** %


\hauteurLargeurCadreReponse		{8mm}{5cm}


% ^^^^^^^^^^^^^^^^^^^^^^^^^^^^^^^^^^^^^^ %
%               Question      d           %
% ^^^^^^^^^^^^^^^^^^^^^^^^^^^^^^^^^^^^^^ %

\begin{enonce}
Donner l'expression de $K^\circ$ en fonction de $K_A$ et $K_e$
\end{enonce}

\reponse{$K^\circ=\dfrac{K_A}{K_e}$}

\begin{corrige}
On a, d'après les questions précédentes : \minusMedskip
\[
K_A=\frac{a(\text{NH}_3)_\text{eq}\times a(\text{H}_3\text{O}^+)_\text{eq}}{a(\text{NH}^+_4)_\text{eq}\times a(\text{H}_2\text{O})_\text{eq}}\qquad \text{et}\qquad K_e=\frac{a(\text{HO}^-)_\text{eq}\times a(\text{H}_3\text{O}^+)_\text{eq}}{a(\text{H}_2\text{O})^2_\text{eq}}.
\]
Donc $\frac{K_A}{K_e}=\frac{a(\text{NH}_3)_\text{eq}\times a(\text{H}_2\text{O})_\text{eq}}{a(\text{NH}^+_4)_\text{eq}\times a(\text{HO}^-)_\text{eq}}$. On en déduit donc que $K^\circ=\dfrac{K_A}{K_e}$.
\end{corrige}

% ************************************** %

% ^^^^^^^^^^^^^^^^^^^^^^^^^^^^^^^^^^^^^^ %
%               Question      e           %
% ^^^^^^^^^^^^^^^^^^^^^^^^^^^^^^^^^^^^^^ %

\begin{enonce}
À \SI{25}{\celsius}, on donne $\mathrm{p}K_A=-\log_{10}(K_A)=\num{9,25}$ et $\mathrm{p}K_e=-\log_{10}(K_e)=14$. 

Calculer $K^\circ$
\end{enonce}

\reponse{$10^{\np{4,75}}$}
%\reponse{\num{5,6e4}}

\begin{corrige}
On a $K^\circ=\frac{K_A}{K_e}=\frac{10^{\np{-9,25}}}{10^{-14}}=10^{\np{4,75}}$.
%On a $K^\circ=\frac{K_A}{K_e}=\frac{10^{-9,25}}{10^{-14}}=10^{4,75}
%=\num{5,6e4}$.
\end{corrige}

% ************************************** %





% =============================================================================== %
\finEntrainement
% =============================================================================== %








% •••••••••••••••••••••••••••••••••••••••••••••••••••••••••••••••••••••••••••• %
% •••••••••••••••••••••••••••••••••••••••••••••••••••••••••••••••••••••••••••• %
\sectionFicheEntrainement{Composition finale d'un système siège d'une réaction chimique}
% •••••••••••••••••••••••••••••••••••••••••••••••••••••••••••••••••••••••••••• %
% •••••••••••••••••••••••••••••••••••••••••••••••••••••••••••••••••••••••••••• %




% ********************************************************************* %
%                           ENTRAÎNEMENT           6                     % 
% ********************************************************************* %
% ================ Métadonnées sur l'entraînement ===================== %

\titreEntrainementFacultatif	{Sens d'évolution d'une réaction}
\hauteurLargeurCadreReponse		{6mm}{12mm}
\dureeResolutionFacultative		{2} % 1, 2, 3 ou 4
\typeEntrainement				{QCM} % Newton ou QCM ou AN ou vide (exercice "normal")
\basiqueEtTransversal			{Y}
\nombreColonnesQuestions		{1} % vide, 1, 2, 3, etc.
\avecPlusieursQuestions			{Y} % Y ou N
\initialisationEntrainement
% ===================================================================== %

On considère la transformation d'équation : \vspace{-.5mm}
\begin{center}
{CH$_3$COOH}$_\text{(aq)}$ + F$^-_\text{(aq)}$ = CH$_3$COO$^-_\text{(aq)}$ + HF$_\text{(aq)}$\vspace{-.5mm}
\end{center}
dont la constante d'équilibre à \SI{25}{\celsius} est $K^\circ=10^{\np{-1,6}}$.

On réalise cette réaction en partant de différentes concentrations initiales de réactifs et de produits. 

\smallskip
Pour chacun des cas ci-dessous, déterminer le sens d'évolution de la réaction.
% =============================================================================== %
\debutEntrainement
% =============================================================================== %

% ^^^^^^^^^^^^^^^^^^^^^^^^^^^^^^^^^^^^^^ %
%               Question     a            %
% ^^^^^^^^^^^^^^^^^^^^^^^^^^^^^^^^^^^^^^ %

\begin{enonce}
$\left[\text{CH}_3\text{COOH}\right]_i=\left[\text{F}^-\right]_i=\SI{1e-1}{\mole\per\liter}$ \qquad et\qquad $\left[\text{CH}_3\text{COO}^-\right]_i=\left[\text{HF}\right]_i=\SI{0}{\mole\per\liter}$
\begin{listeQCM2Colonnes}
	\item sens direct
	\item sens indirect
	\item pas d'évolution
\end{listeQCM2Colonnes}
\end{enonce}

\reponse{\reponseA}

\begin{corrige}
À l'état initial, $Q_i=\frac{\left[\text{HF}\right]_i\times\left[\text{CH}_3\text{COO}^-\right]_i}{\left[\text{CH}_3\text{COOH}\right]_i\times\left[\text{F}^-\right]_i}=0 <K^\circ$. La réaction évolue donc dans le sens direct.
\end{corrige}

% ************************************** %


% ^^^^^^^^^^^^^^^^^^^^^^^^^^^^^^^^^^^^^^ %
%               Question       b          %
% ^^^^^^^^^^^^^^^^^^^^^^^^^^^^^^^^^^^^^^ %

\begin{enonce}
$\left[\text{CH}_3\text{COOH}\right]_i=
\left[\text{F}^-\right]_i=
\left[\text{CH}_3\text{COO}^-\right]_i=\SI{1e-1}{\mole\per\liter}$ \qquad et\qquad 
$\left[\text{HF}\right]_i=\SI{0}{\mole\per\liter}$
\begin{listeQCM2Colonnes}
	\item sens direct
	\item sens indirect
	\item pas d'évolution
\end{listeQCM2Colonnes}
\end{enonce}

\reponse{\reponseA}

\begin{corrige}
À l'état initial, $Q_i=\frac{\left[\text{HF}\right]_i\times\left[\text{CH}_3\text{COO}^-\right]_i}{\left[\text{CH}_3\text{COOH}\right]_i\times\left[\text{F}^-\right]_i}=0<K^\circ$. La réaction évolue donc dans le sens direct.
\end{corrige}

% ************************************** %


% ^^^^^^^^^^^^^^^^^^^^^^^^^^^^^^^^^^^^^^ %
%               Question       c          %
% ^^^^^^^^^^^^^^^^^^^^^^^^^^^^^^^^^^^^^^ %

\begin{enonce}
$\left[\text{CH}_3\text{COOH}\right]_i=
\left[\text{F}^-\right]_i = 
\left[\text{CH}_3\text{COO}^-\right]_i=
\left[\text{HF}\right]_i=\SI{1,0e-1}{\mole\per\liter}
$
\begin{listeQCM2Colonnes}
	\item sens direct
	\item sens indirect
	\item pas d'évolution
\end{listeQCM2Colonnes}
\end{enonce}


\reponse{\reponseB}

\begin{corrige}
À l'état initial, $Q_i=\frac{\left[\text{HF}\right]_i\times\left[\text{CH}_3\text{COO}^-\right]_i}{\left[\text{CH}_3\text{COOH}\right]_i\times\left[\text{F}^-\right]_i}=1>K^\circ$. La réaction évolue donc dans le sens indirect.
\end{corrige}

% ************************************** %


% ^^^^^^^^^^^^^^^^^^^^^^^^^^^^^^^^^^^^^^ %
%               Question       d          %
% ^^^^^^^^^^^^^^^^^^^^^^^^^^^^^^^^^^^^^^ %

\begin{enonce}
$\left[\text{CH}_3\text{COOH}\right]_i=\SI{8,0e-4}{\mole\per\liter}$  \qquad et\qquad  $\left[\text{F}^-\right]_i=
\left[\text{HF}\right]_i=\SI{4,0e-3}{\mole\per\liter}$ \\ 
et $\left[\text{CH}_3\text{COO}^-\right]_i=\SI{2,0e-5}{\mole\per\liter}$ 

\begin{listeQCM2Colonnes}
	\item sens direct
	\item sens indirect
	\item pas d'évolution
\end{listeQCM2Colonnes}
\end{enonce}

\reponse{\reponseC}

\begin{corrige}
À l'état initial, $Q_i=\frac{\left[\text{HF}\right]_i\times\left[\text{CH}_3\text{COO}^-\right]_i}{\left[\text{CH}_3\text{COOH}\right]_i\times\left[\text{F}^-\right]_i}=\num{2,5e-2}=10^{-1,6}=K^\circ$. Ainsi, le système est à l'équilibre et n'évolue pas.
\end{corrige}

% ************************************** %


% =============================================================================== %
\finEntrainement
% =============================================================================== %



\newpage



% ********************************************************************* %
%                           ENTRAÎNEMENT            7                    % 
% ********************************************************************* %
% ================ Métadonnées sur l'entraînement ===================== %

\titreEntrainementFacultatif	{Détermination du réactif limitant}
\hauteurLargeurCadreReponse		{6mm}{12mm}
\dureeResolutionFacultative		{1} % 1, 2, 3 ou 4
\typeEntrainement				{QCM} % Newton ou QCM ou AN ou vide (exercice "normal")
\nombreColonnesQuestions		{1} % vide, 1, 2, 3, etc.
\avecPlusieursQuestions			{N} % Y ou N
\initialisationEntrainement
% ===================================================================== %

%Une solution de chlorure de fer (III) (Fe$^{3+}$ ; 3 Cl$^-$ ) réagit avec de l'hydroxyde de sodium  (Na$^+$ ; HO$^-$ ) en formant un précipité d'hydroxyde de fer {Fe(OH)$_3$}$_\text{(s)}$, aussi connu sous le nom de rouille. On donne : 

On considère la réaction entre les ions fer (III) et les ions hydroxyde, formant un précipité d'hydroxyde de fer {Fe(OH)$_3$}$_\text{(s)}$, aussi connu sous le nom de rouille. L'équation de la réaction est :\vspace{-1.5mm}
\begin{center}
Fe$^{3+}_\text{(aq)}$ + 3 HO$^-_\text{(aq)}$ = {Fe(OH)$_3$}$_\text{(s)}$\vspace{-1.5mm}
\end{center}
À l'instant initial, on mélange une solution de chlorure de fer (III) (Fe$^{3+}$ ; 3 Cl$^-$) avec une solution de soude (hydroxyde de sodium (Na$^+$ ; HO$^-$)) de sorte à obtenir les conditions suivantes : \vspace{-2mm}
\begin{center}
\renewcommand{\arraystretch}{1.2}
\begin{tabular}{|l|c|c|c|c|}
\hline
& Fe$^{2+}$ & Cl$^-$ & Na$^+$ & HO$^-$\\
\hline
\textbf{Quantité de matière initiale}& \SI{3,0e-2}{\mole} & \SI{9,0e-2}{\mole} & \SI{6,0e-2}{\mole} & \SI{6,0e-2}{\mole}\\
\hline
\end{tabular}
\end{center}
Déterminer le réactif limitant.

%\begin{center} 
%\renewcommand{\arraystretch}{1.5}
%\begin{tabular}{|l|c|c||c|}
%\hline
%&\multicolumn{1}{c@{\protect\makebox[0cm]{$+$}}}{~~~~$\mathrm{Fe^{3+}}_\text{(aq)}$~~~~}
%&\multicolumn{1}{c@{\protect\makebox[0cm]{$=$}}}{~~~~$3\,\mathrm{HO^-}_\text{(aq)}$~~~~}
%&\multicolumn{1}{c|}{~~~~$\mathrm{Fe(OH)_3}_\text{(s)}$~~~}
%\\
%\hline
%\textbf{\'Etat initial}	& \SI{3,0e-2}{\mole}		&	\SI{6,0e-2}{\mole}	& 0
%\\ \hline
%\end{tabular}
%\end{center}

% =============================================================================== %
\debutEntrainement
% =============================================================================== %

% ^^^^^^^^^^^^^^^^^^^^^^^^^^^^^^^^^^^^^^ %
%               Question                 %
% ^^^^^^^^^^^^^^^^^^^^^^^^^^^^^^^^^^^^^^ %

\begin{enonce}
\begin{listeQCM4Colonnes}
	\item $\mathrm{Fe^{3+}}_\text{(aq)}$
	\item $\mathrm{HO^-}_\text{(aq)}$
	\item $\mathrm{Fe(OH)_3}_\text{(s)}$
	\item Il n'y en a pas
\end{listeQCM4Colonnes}
\bigskip
\end{enonce}

\reponse{\reponseB}
\begin{corrige}
On calcule, pour chaque réactif, le rapport entre sa quantité de matière initiale et son nombre stœchiométrique. Le réactif pour lequel ce rapport est le plus faible est le réactif limitant. 

On trouve
$\frac{n(\text{Fe}^{3+})_i}{1}$ = \SI{3,0e-2}{\mole} et $\frac{n(\text{OH}^-)_i}{3}$ = \SI{2,0e-2}{\mole}. 

L'ion hydroxyde HO$^-$ est donc le réactif limitant.

{\itshape Remarque : on ne prend pas en compte les ions $\mathrm{Na}^+$ ou $\mathrm{Cl}^-$ car ce sont des ions spectateurs et non des réactifs.}
\end{corrige}

% ************************************** %




% =============================================================================== %
\finEntrainement
% =============================================================================== %






% ********************************************************************* %
%                           ENTRAÎNEMENT            8                    % 
% ********************************************************************* %
% ================ Métadonnées sur l'entraînement ===================== %

\titreEntrainementFacultatif	{Transformation totale}
\hauteurLargeurCadreReponse		{6mm}{25mm}
\dureeResolutionFacultative		{2} % 1, 2, 3 ou 4
\typeEntrainement				{AN} % Newton ou QCM ou AN ou vide (exercice "normal")
\nombreColonnesQuestions		{1} % vide, 1, 2, 3, etc.
\avecPlusieursQuestions			{Y} % Y ou N
\initialisationEntrainement
% ===================================================================== %

On considère la réaction de combustion du butane à l'état gazeux suivante, ainsi que les concentrations initiales des réactifs :   \vspace{-4mm}
\begin{center}
{2 C$_4$H$_{10}$}$_\text{(g)}$ + 13 {O$_2$}$_\text{(g)}$ $\longrightarrow$ 8 {CO$_2$}$_\text{(g)}$ + 10 H$_2$O$_\text{(g)}$

%\end{center}
%On donne : 
%\begin{center}
\renewcommand{\arraystretch}{1.2}
\begin{tabular}{|l|c|c|c|c|}
\hline
& C$_4$H$_{10}$ & O$_2$ & CO$_2$ & H$_2$O \\
\hline
\textbf{Quantité de matière initiale}& $n_1=\SI{0,10}{\mol}$ & $n_2=\SI{0,65}{\mol}$ & $\SI{0}{\mol}$ & $\SI{0}{\mol}$\\
\hline
\end{tabular}
\vspace{-1mm}
\end{center}
Sachant que la réaction est totale, déterminer : 

%\begin{center}
%\renewcommand{\arraystretch}{1.5}
%\begin{tabular}{|l|c|c||c|c|}
%\hline
%&\multicolumn{1}{c@{\protect\makebox[0cm]{$+$}}}{~~~~$2\,\mathrm{C_4H_{10}}_\text{(g)}$~~~~}
%&\multicolumn{1}{c@{\protect\makebox[0cm]{$\longrightarrow$}}}{~~~~$13\,\mathrm{O_2}_\text{(g)}$~~~~}
%&\multicolumn{1}{c@{\protect\makebox[0cm]{$+$}}}{~~~~$8\,\mathrm{CO_2}_\text{(g)}$~~~~}
%&\multicolumn{1}{c|}{~~~~$10\,\mathrm{H_2O}_\text{(g)}$~~~~}
%\\
%\hline
%\textbf{\'Etat initial}		& 
%$n_1=\SI{0,10}{\mol}$ 		&
%$n_2=\SI{0,65}{\mol}$		&
%0							&
%0
%\\ \hline
%\end{tabular}
%\end{center}


% =============================================================================== %
\debutEntrainement
% =============================================================================== %

% ^^^^^^^^^^^^^^^^^^^^^^^^^^^^^^^^^^^^^^ %
%               Question      a           %
% ^^^^^^^^^^^^^^^^^^^^^^^^^^^^^^^^^^^^^^ %

\begin{enonce}
L'avancement maximal $\xi_\text{max}$ pour cette transformation
\end{enonce}

\reponse{$\SI{5,0e-2}{\mol}$}

\begin{corrige}
On a $\frac{n_1}{2}=\frac{n_2}{13}=\num{5,0e-2}\si{\mole}$ : les réactifs ont donc été introduits en proportions stœchiométriques. Dans ce cas, il n'y a pas de réactif limitant (ou alors tous les réactifs sont limitants). 

L'avancement maximal est alors $\xi_\text{max}=\num{5,0e-2}\si{\mole}$.
\end{corrige}

% ************************************** %

% ^^^^^^^^^^^^^^^^^^^^^^^^^^^^^^^^^^^^^^ %
%               Question      b           %
% ^^^^^^^^^^^^^^^^^^^^^^^^^^^^^^^^^^^^^^ %

\begin{enonce}
La quantité de matière de dioxyde de carbone (CO$_2$) à l'état final
\end{enonce}

\reponse{$\SI{4,0e-1}{\mol}$}

\begin{corrige}
On écrit un tableau d'avancement pour la réaction totale : \minusMedskip
\begin{center}
\renewcommand{\arraystretch}{1.2}
\begin{tabular}{|l|c|c||c|c|}
\hline
&\multicolumn{1}{c@{\protect\makebox[0cm]{$+$}}}{~~~~$2\,\mathrm{C_4H_{10}}_\text{(g)}$~~~~}
&\multicolumn{1}{c@{\protect\makebox[0cm]{$\longrightarrow$}}}{~~~~$13\,\mathrm{O_2}_\text{(g)}$~~~~}
&\multicolumn{1}{c@{\protect\makebox[0cm]{$+$}}}{~~~~$8\,\mathrm{CO_2}_\text{(g)}$~~~~}
&\multicolumn{1}{c|}{~~~~$10\,\mathrm{H_2O}_\text{(g)}$~~~~}
\\
\hline
\textbf{\'Etat initial}		& 
$n_1$				 		&
$n_2$						&
0							&
0							\\ 
\hline
\textbf{\'Etat final} 		&
$n_1-2\xi_\text{max}$		&
$n_2-13\xi_\text{max}$		&
$8\xi_\text{max}$			&
$10\xi_\text{max}$			\\ 
\hline
\end{tabular}
\minusMedskip
\end{center}
Comme la réaction est totale, l'avancement atteint à l'état final correspond à l'avancement maximal $\xi_\text{max}$ calculé à la question précédente. On a donc $n(\text{CO}_2)_f=8\xi_\text{max}=\num{4,0e-1}\si{\mole}$.
\end{corrige}

% ************************************** %

% =============================================================================== %
\finEntrainement
% =============================================================================== %







% ********************************************************************* %
%                           ENTRAÎNEMENT        9                        % 
% ********************************************************************* %
% ================ Métadonnées sur l'entraînement ===================== %

\titreEntrainementFacultatif	{Une autre transformation totale}
\hauteurLargeurCadreReponse		{6mm}{12mm}
\dureeResolutionFacultative		{2} % 1, 2, 3 ou 4
\typeEntrainement				{QCM} % Newton ou QCM ou AN ou vide (exercice "normal")
\nombreColonnesQuestions		{1} % vide, 1, 2, 3, etc.
\avecPlusieursQuestions			{Y} % Y ou N
\initialisationEntrainement
% ===================================================================== %

% =============================================================================== %
\debutEntrainement
% =============================================================================== %

On s'intéresse à la réaction des ions argent avec le cuivre selon l'équation de réaction : 
\[
2~\text{Ag}^+_\text{(aq)}+\text{Cu}_\text{(s)}\longrightarrow \text{Cu}^{2+}_\text{(aq)}+2~\text{Ag}_\text{(s)}.
\]
Cette réaction est totale. On mélange initialement un volume $V=\SI{20}{m\liter}$ d'une solution contenant des ions argent ($\text{Ag}^+$) à la concentration $C=\SI{0,25}{\mol\per\liter}$ avec une masse $m=\SI{0,254}{\gram}$ de cuivre solide ($\text{Cu}$). 

On donne la masse molaire du cuivre  $M_{\text{Cu}}= \SI{63,5}{\gram\per\mol}$ et celle de l'argent $M_{\text{Ag}}=\SI{107}{\gram\per\mol}$.

% ^^^^^^^^^^^^^^^^^^^^^^^^^^^^^^^^^^^^^^ %
%               Question     a            %
% ^^^^^^^^^^^^^^^^^^^^^^^^^^^^^^^^^^^^^^ %

\begin{enonce}
Quel est le réactif limitant ?
\begin{listeQCM3Colonnes} %proposition de mise en page
	\item $\text{Ag}^+_\text{(aq)}$
	\item $\mathrm{Cu}_\text{(s)}$
	\item Il n'y en a pas
\end{listeQCM3Colonnes}
%\begin{listeQCM1Colonne}
%	\item $\text{Ag}^+_\text{(aq)}$
%	\item $\mathrm{Cu}_\text{(s)}$
%	\item Il n'y en a pas
%\end{listeQCM1Colonne}
\end{enonce}

\reponse{\reponseA}

\begin{corrige}
On calcule dans un premier temps les quantités de matière initiales de tous les réactifs : \minusSmallskip
\begin{align*}
 n(\text{Ag}^+)_i&=n_1=C\times V=\SI{0,25}{\mol\per\liter} \times \SI{20e-3}{\liter}=\SI{5,0e-3}{\mole}\\
 n(\text{Cu})_i&=n_2=\frac{m}{M_{\text{Cu}}}=\frac{\SI{0,254}{\gram}}{\SI{63,5}{\gram\per\mol}}=\SI{4,0e-3}{\mole}.\minusSmallskip
\end{align*}
On calcule ensuite les rapports entre les quantités de matière initiales et les nombres stœchiométriques : 
\[
\frac{ n(\text{Ag}^+)_i}{2}=\SI{2,5e-3}{\mole} <\frac{n(\text{Cu})_i}{1}=\SI{4,0e-3}{\mole}
\]
Le réactif limitant est donc Ag$^+$.
\end{corrige}

% ************************************** %

% ^^^^^^^^^^^^^^^^^^^^^^^^^^^^^^^^^^^^^^ %
%               Question     b            %
% ^^^^^^^^^^^^^^^^^^^^^^^^^^^^^^^^^^^^^^ %

\begin{enonce}
À la fin de la réaction, la quantité de matière de $\mathrm{Cu}_\text{(s)}$ vaut :
	\begin{listeQCM3Colonnes}%proposition demise en page
	\item \SI{1.5}{\milli\mol}
	\item \SI{2.5}{\milli\mol}
	\item \SI{0}{\milli\mol}
	\end{listeQCM3Colonnes}		
%	\begin{listeQCM1Colonne}
%	\item \SI{1.5}{\milli\mol}
%	\item \SI{2.5}{\milli\mol}
%	\item \SI{0}{\milli\mol}
%	\end{listeQCM1Colonne}		
\end{enonce}

\reponse{\reponseA{}}

\begin{corrige}
On dresse un tableau d'avancement pour cette réaction : \minusMedskip
\begin{center}%nouvelle version
\renewcommand{\arraystretch}{1.2}
\begin{tabular}{|l|c|c||c|c|}
  \hline
&
%
\multicolumn{1}{c@{\protect\makebox[0cm]{$\mathrm{+}$}}}{%
~~~~$\mathrm{2~{Ag^+}_\text{(aq)}}$~~~~}                   &
\multicolumn{1}{c@{\protect\makebox[0cm]{$=$}}}{%
~~~~$\mathrm{{Cu}_\text{(s)}}$~~~~}                        &
\multicolumn{1}{c@{\protect\makebox[0cm]{$\mathrm{+}$}}}{%
~~~~$\mathrm{{Cu^{2+}}_\text{(aq)}}$~~~~}                        &
\multicolumn{1}{c|}{%
~~~~$\mathrm{2~{Ag}_\text{(s)}}$~~~}              \\ \hline
%
\textbf{\'Etat initial}			&
$n_1$                       	&
$n_2$							&
0								&
0      							\\ \hline
%
\textbf{\'Etat final}			& 
$n_1-2\xi_\text{max}$			&
$n_2-\xi_\text{max}$			&
$\xi_\text{max}$                &
$2\xi_\text{max}$				\\ \hline
\end{tabular}
\minusMedskip
\end{center}
La réaction est totale, donc l'avancement final est égal à l'avancement maximal.\\
Le réactif limitant est l'ion argent (Ag$^+$), donc l'avancement final est $\xi_\text{max}=\frac{n_1}{2}=\SI{2,5}{m\mole}$. \\
À l'état final, on a donc $n(\text{Cu})_f=\SI{4.0}{\milli\mol}-\SI{2.5}{\milli\mol}=\SI{1,5}{\milli\mol}$.
%\[
%n(\text{Ag}^+)_f=\SI{0}{m\mole}~~;~~n(\text{Cu})_f=\SI{1,5}{m\mole}~~;~~n(\text{Cu}^{2+})_f=\SI{2,5}{m\mole}~~;~~n(\textrm{Ag})_f=\SI{5}{m\mole}
%\]
\end{corrige}

% ************************************** %

% =============================================================================== %
\finEntrainement
% =============================================================================== %



\newpage



% ********************************************************************* %
%                           ENTRAÎNEMENT            10                    % 
% ********************************************************************* %
% ================ Métadonnées sur l'entraînement ===================== %

\titreEntrainementFacultatif	{Loi d'action de masse et composition à l'équilibre}
\hauteurLargeurCadreReponse		{12mm}{70mm}
\dureeResolutionFacultative		{3} % 1, 2, 3 ou 4
\typeEntrainement				{Newton} % Newton ou QCM ou AN ou vide (exercice "normal")
\nombreColonnesQuestions		{1} % vide, 1, 2, 3, etc.
\avecPlusieursQuestions			{Y} % Y ou N
\initialisationEntrainement
% ===================================================================== %

%On considère la transformation d'équation :
%\begin{center}
%{PhCOO$^-$}$_\text{(aq)}$ + {H$_3$O$^+$}$_\text{(aq)}$ = {PhCOOH}$_\text{(s)}$ + {H$_2$O}$_\text{(l)}$
%\end{center}
À l'instant initial, on mélange un volume $V_1$ d'une solution aqueuse d'ions benzoate (PhCOO$^-$) à la concentration $C_1$ et un volume $V_2$ d'une solution aqueuse d'ions oxonium (H$_3$O$^+$) à la concentration $C_2$.

On donne l'équation de la réaction et son tableau d'avancement en quantité de matière :

\begin{center}
\renewcommand{\arraystretch}{1.5}
\begin{tabular}{|l|c|c||c|c|}
  \hline
&
%
\multicolumn{1}{c@{\protect\makebox[0cm]{$\mathrm{+}$}}}{%
~~~~$\mathrm{PhCOO^-_\text{(aq)}}$~~~~}                   &
\multicolumn{1}{c@{\protect\makebox[0cm]{$=$}}}{%
~~~~$\mathrm{H_3O^+_\text{(aq)}}$~~~~}                        &
\multicolumn{1}{c@{\protect\makebox[0cm]{$\mathrm{+}$}}}{%
~~~~$\mathrm{PhCOOH_\text{(s)}}$~~~~}                        &
\multicolumn{1}{c|}{%
~~~~$\mathrm{H_2O_\text{($\ell$)}}$~~~}              \\ \hline
%
\textbf{\'Etat initial}		&
$C_1V_1$					&
$C_2V_2$					&
0							&
\text{excès}				\\ \hline
%
\textbf{\'Etat final} 		&
$C_1V_1-\xi$				&
$C_2V_2-\xi$				&
$\xi$						&
\text{excès}				\\ \hline
\end{tabular}
\end{center}

% =============================================================================== %
\debutEntrainement
% =============================================================================== %

% ^^^^^^^^^^^^^^^^^^^^^^^^^^^^^^^^^^^^^^ %
%               Question       a          %
% ^^^^^^^^^^^^^^^^^^^^^^^^^^^^^^^^^^^^^^ %

\begin{enonce}
À l'aide de la loi d'action de masse, exprimer la constante d'équilibre $K^\circ$ associée à cette réaction, en fonction de $C_1$, $C_2$, $V_1$, $V_2$, $C^\circ$ et $\xi$. \minusBigskip

\end{enonce}

\reponse{$\frac{\left(C^\circ(V_1+V_2)\right)^2}{\left(C_1V_1-\xi\right)\times\left(C_2V_2-\xi\right)}$}

\begin{corrige}
D'après la loi d'action de masse, $K^\circ=\frac{a(\text{PhCOOH})_\text{eq}\times a(\text{H}_2\text{O})_\text{eq}}{a(\text{PhCOO}^-)_\text{eq}\times a(\text{H}_3\text{O}^+)_\text{eq}}$. Comme PhCOOH est un solide seul dans sa phase, son activité vaut 1. Comme H$_2$O est le solvant, son activité vaut 1. L'activité des espèces aqueuses s'exprime en fonction de leur concentration et de $C^\circ$. 

Avec les expressions du tableau d'avancement, on a alors : 
\[
K^\circ=\frac{1\times 1}{\left(\frac{1}{C^\circ}\frac{C_1V_1-\xi}{V_1+V_2}\right)\times\left(\frac{1}{C^\circ}\frac{C_2V_2-\xi}{V_1+V_2}\right)}=
\frac{\left(C^\circ(V_1+V_2)\right)^2}{\left(C_1V_1-\xi\right)\times\left(C_2V_2-\xi\right)}.
\]
\end{corrige}

% ************************************** %

% ^^^^^^^^^^^^^^^^^^^^^^^^^^^^^^^^^^^^^^ %
%               Question      b           %
% ^^^^^^^^^^^^^^^^^^^^^^^^^^^^^^^^^^^^^^ %

\begin{enonce}
En déduire l'équation du second degré permettant de déterminer la valeur de $\xi$.

\end{enonce}

\reponse{$\xi^2-\xi(C_1V_1+C_2V_2)+C_1C_2V_1V_2-\frac{\left[C^\circ(V_1+V_2)\right]^2}{K^\circ}=0$}

\begin{corrige}
À partir de la relation précédente on déduit $\left(C_1V_1-\xi\right)\times\left(C_2V_2-\xi\right)=\frac{\left(C^\circ(V_1+V_2)\right)^2}{K^\circ}$.

Après développement on obtient \minusBigskip
\[
\xi^2-\xi(C_1V_1+C_2V_2)+C_1C_2V_1V_2-\frac{\left(C^\circ(V_1+V_2)\right)^2}{K^\circ}=0.
\]
Au passage, la formule obtenue est bien homogène : chaque terme est homogène à une quantité de matière au carré.
\end{corrige}

% ************************************** %




% =============================================================================== %
\finEntrainement
% =============================================================================== %





% ********************************************************************* %
%                           ENTRAÎNEMENT              11                  % 
% ********************************************************************* %
% ================ Métadonnées sur l'entraînement ===================== %

\titreEntrainementFacultatif	{À la recherche de l'équilibre}
\hauteurLargeurCadreReponse		{12mm}{8cm} % La largeur de la zone texte fait bien 16cm ?
\dureeResolutionFacultative		{2} % 1, 2, 3 ou 4
\typeEntrainement				{Newton} % Newton ou QCM ou AN ou vide (exercice "normal")
\basiqueEtTransversal			{Y}
\nombreColonnesQuestions		{1} % vide, 1, 2, 3, etc.
\avecPlusieursQuestions			{Y} % Y ou N
\initialisationEntrainement
% ===================================================================== %

La loi d'action de masse permet de déterminer l'avancement $\xi$ ou l'avancement volumique $\xi_v$ à l'équilibre. 

Mettre ces différentes lois d'action de masse sous la forme d'une équation du second degré en $\xi$ ou $\xi_v$.

\bigskip


% =============================================================================== %
\debutEntrainement
% =============================================================================== %


% ^^^^^^^^^^^^^^^^^^^^^^^^^^^^^^^^^^^^^^ %
%               Question   b              %
% ^^^^^^^^^^^^^^^^^^^^^^^^^^^^^^^^^^^^^^ %

\begin{enonce}
$K^\circ=\frac{\xi^2_v}{(C_1-\xi_v)\times(C_2-\xi_v)}$
\end{enonce}

\reponse{$\xi^2_v(1-K^\circ)+\xi_v K^\circ\left(C_1+C_2\right)-K^\circ C_1C_2=0$}

%\begin{corrige}
%En ré-arrangeant l'expression donnée, on obtient : 
%\[
%\xi^2_v(1-K^\circ)+\xi_v K^\circ\left(C_1+C_2\right)-K^\circ C_1C_2=0
%\]
%\end{corrige}

% ************************************** %


% ^^^^^^^^^^^^^^^^^^^^^^^^^^^^^^^^^^^^^^ %
%               Question  c               %
% ^^^^^^^^^^^^^^^^^^^^^^^^^^^^^^^^^^^^^^ %

\begin{enonce}
$K^\circ=\frac{\xi_v\left(C_2+\xi_v\right)}{(C_1-\xi_v)\times C^\circ}$
\end{enonce}

\reponse{$\xi^2_v+\xi_v\left(C_2+K^\circ C^\circ\right)-K^\circ C_1C^\circ=0$}

%\begin{corrige}
%En ré-arrangeant l'expression donnée, on obtient : 
%\[
%\xi^2_v+\xi_v\left(C_2+K^\circ C^\circ\right)-K^\circ C_1C^\circ=0
%\]
%\end{corrige}

% ************************************** %


% ^^^^^^^^^^^^^^^^^^^^^^^^^^^^^^^^^^^^^^ %
%               Question  d               %
% ^^^^^^^^^^^^^^^^^^^^^^^^^^^^^^^^^^^^^^ %

\begin{enonce}
$K^\circ=\frac{\left(\frac{\xi RT}{V}\right)^2}{\left(\frac{(n_1-\xi)RT}{V}\right)\times \left(\frac{(n_2-\xi)RT}{V}\right)}$
\end{enonce}

\reponse{$\xi^2(K^\circ-1)-\xi K^\circ(n_1+n_2)+K^\circ n_1n_2=0$}

%\begin{corrige}
%En ré-arrangeant l'expression donnée, on obtient : 
%\[
%\xi^2(K^\circ-1)-\xi K^\circ(n_1+n_2)+K^\circ n_1n_2=0
%\]
%\end{corrige}

% ************************************** %

% ^^^^^^^^^^^^^^^^^^^^^^^^^^^^^^^^^^^^^^ %
%               Question    e             %
% ^^^^^^^^^^^^^^^^^^^^^^^^^^^^^^^^^^^^^^ %

\begin{enonce}
$K^\circ=\frac{\left(\frac{\xi RT}{V}\right)\cdot P^\circ}{\left(\frac{(n-2\xi)RT}{V}\right)^2}$
\end{enonce}

\reponse{$4K^\circ\xi^2-\xi\left(4K^\circ n+\frac{P^\circ V}{RT}\right)+K^\circ n^2=0$}

%\begin{corrige}
%En ré-arrangeant l'expression donnée, on obtient : 
%\[
%4K^\circ\xi^2-\xi\left(4K^\circ n+\frac{P^\circ V}{RT}\right)+K^\circ n^2=0
%\]
%\end{corrige}

% ************************************** %


% ^^^^^^^^^^^^^^^^^^^^^^^^^^^^^^^^^^^^^^ %
%               Question  f               %
% ^^^^^^^^^^^^^^^^^^^^^^^^^^^^^^^^^^^^^^ %

\begin{enonce}
$K^\circ=\frac{\left(\frac{\xi}{n-\xi}\cdot P\right)P^\circ}{\left(\frac{(n-2\xi)}{n-\xi}\cdot P\right)^2}$
\end{enonce}

\reponse{$\xi^2(4K^\circ P+P^\circ)-\xi(4n K^\circ P+n P^\circ)+K^\circ n^2 P=0 $}

%\begin{corrige}
%En ré-arrangeant l'expression donnée, on obtient : 
%\[
%\xi^2(4K^\circ P+P^\circ)-\xi(4n K^\circ P+n P^\circ)+K^\circ n^2 P=0
%\]
%\end{corrige}

% ************************************** %



% =============================================================================== %
\finEntrainement
% =============================================================================== %



\newpage



% ********************************************************************* %
%                           ENTRAÎNEMENT              12                  % 
% ********************************************************************* %
% ================ Métadonnées sur l'entraînement ===================== %

\titreEntrainementFacultatif	{Calcul de l'avancement à l'équilibre}
\hauteurLargeurCadreReponse		{6mm}{18mm}
\dureeResolutionFacultative		{3} % 1, 2, 3 ou 4
\typeEntrainement				{AN} % Newton ou QCM ou AN ou vide (exercice "normal")
\basiqueEtTransversal			{Y}
\nombreColonnesQuestions		{1} % vide, 1, 2, 3, etc.
\avecPlusieursQuestions			{Y} % Y ou N
\initialisationEntrainement
% ===================================================================== %



Dans chacune des situations suivantes, une réaction se produit dans le sens direct. On indique que son  l'avancement maximal est $\xi_\text{v,max}=\SI{1,0e-1}{\mole\per\liter}$. 

La loi d'action de masse donne l'équation dont est solution l'avancement volumique $\xi_v$. 

\medskip

Calculer $\xi_v$.

% =============================================================================== %
\debutEntrainement
% =============================================================================== %

% ^^^^^^^^^^^^^^^^^^^^^^^^^^^^^^^^^^^^^^ %
%               Question      a           %
% ^^^^^^^^^^^^^^^^^^^^^^^^^^^^^^^^^^^^^^ %

\begin{enonce}
$\xi^2_v(1-K^\circ)+\xi_v K^\circ\left(C_1+C_2\right)-K^\circ C_1C_2=0
\quad\text{avec}\quad
\begin{cases}
K^\circ=\num{2,0}\\
C_2=2C_1=\SI{1,0e-1}{\mole\per\liter}
\end{cases}
$
\end{enonce}
\reponse{$\SI{7,6e-2}{\mole\per\liter} $}
\begin{corrige}
La résolution du polynôme du second degré donne deux solutions : \vspace{-2mm}
\[
\xi_{v,1}=\SI{7,6e-2}{\mole\per\liter}\quad\text{et}\quad\xi_{v,2}=\SI{5,2e-1}{\mole\per\liter}. \vspace{-1mm}
\]
L'avancement final ne peut pas être supérieur à l'avancement maximal $\xi_\text{v,max}=\SI{1,0e-1}{\mole\per\liter}$. On en déduit donc que $\xi_v=\SI{7,6e-2}{\mole\per\liter}$.
\end{corrige}

% ************************************** %

% ^^^^^^^^^^^^^^^^^^^^^^^^^^^^^^^^^^^^^^ %
%               Question       b          %
% ^^^^^^^^^^^^^^^^^^^^^^^^^^^^^^^^^^^^^^ %

\begin{enonce}
$\xi^2_v+\xi_v K^\circ C^\circ-K^\circ C_1 C^\circ=0
\quad\text{avec}\quad
\begin{cases}
K^\circ=10^{\np{-1,7}}\\
C_1=\SI{1,0e-1}{\mole\per\liter}
\end{cases}$
\end{enonce}
\reponse{$\SI{3,6e-2}{\mole\per\liter} $}
\begin{corrige}
La résolution du polynôme du second degré donne deux solutions : \vspace{-2mm}
\[
\xi_{v,1}=\SI{3,6e-2}{\mole\per\liter}\quad\text{et}\quad\xi_{v,2}=\SI{-5,6e-2}{\mole\per\liter}. \vspace{-1mm}
\]
Il est indiqué que la réaction se déroule dans le sens direct, donc l'avancement doit être positif. La solution $\xi_{v,2}$ est par conséquent impossible.
On a donc $\xi_v=\SI{3,6e-2}{\mole\per\liter}$, qui est bien inférieur à l'avancement maximal.
\end{corrige}

% ************************************** %


% =============================================================================== %
\finEntrainement
% =============================================================================== %







% •••••••••••••••••••••••••••••••••••••••••••••••••••••••••••••••••••••••••••• %
% •••••••••••••••••••••••••••••••••••••••••••••••••••••••••••••••••••••••••••• %
\sectionFicheEntrainement{Autour des réactions acido-basiques}
% •••••••••••••••••••••••••••••••••••••••••••••••••••••••••••••••••••••••••••• %
% •••••••••••••••••••••••••••••••••••••••••••••••••••••••••••••••••••••••••••• %



% ********************************************************************* %
%                           ENTRAÎNEMENT              13                  % 
% ********************************************************************* %
% ================ Métadonnées sur l'entraînement ===================== %

\titreEntrainementFacultatif	{pH d'une solution}
\hauteurLargeurCadreReponse		{6mm}{12mm}
\dureeResolutionFacultative		{2} % 1, 2, 3 ou 4
\typeEntrainement				{} % Newton ou QCM ou AN ou vide (exercice "normal")
\basiqueEtTransversal			{Y}
\nombreColonnesQuestions		{1} % vide, 1, 2, 3, etc.
\avecPlusieursQuestions			{Y} % Y ou N
\initialisationEntrainement
% ===================================================================== %

La constante d'autoprotolyse de l'eau $K_e=\dfrac{a(\mathrm{HO}^-)\times a(\mathrm{H_3O^+})}{a(\mathrm{H_2O})^2}$ vaut $K_e=10^{-14}$ à \SI{25}{\celsius}.
\medskip


Calculer le pH de la solution dans les cas suivants.

% =============================================================================== %
\debutEntrainement
% =============================================================================== %

% ^^^^^^^^^^^^^^^^^^^^^^^^^^^^^^^^^^^^^^ %
%               Question    a             %
% ^^^^^^^^^^^^^^^^^^^^^^^^^^^^^^^^^^^^^^ %

\begin{enonce}
Une solution telle que $[\mathrm{H_3O^+}]=\SI{5,0e-2}{\mole\per\liter}$
\end{enonce}

\reponse{$\num{1,3}$}

\begin{corrige}
Par définition, $\text{pH}=-\log_{10}\left(a(\text{H}_3\text{O}^+)\right)$. \minusSmallskip

En solution aqueuse diluée, l'activité de H$_3$O$^+$ est $a(\text{H}_3\text{O}^+)=\frac{[\text{H}_3\text{O}^+]}{C^\circ}$. L'expression précédente devient donc : \vspace{-1mm}
\[
\text{pH}=
-\log_{10}\left(
\frac{[\text{H}_3\text{O}^+]}{C^\circ}
\right)=
-\log_{10}\left(
\frac{\SI{5.0E-2}{\mol\per\liter}}{\SI{1,0}{\mol\per\liter}}
\right)=
\num{1,3}.
\]
\end{corrige}

% ************************************** %


% ^^^^^^^^^^^^^^^^^^^^^^^^^^^^^^^^^^^^^^ %
%               Question    b             %
% ^^^^^^^^^^^^^^^^^^^^^^^^^^^^^^^^^^^^^^ %

\begin{enonce}
Une solution telle que $[\mathrm{HO^-}]=\SI{1,0e-2}{\mole\per\liter}$
\end{enonce}

\reponse{$\num{12}$}

\begin{corrige}
Les concentrations $[\text{HO}^-]$ et $[\text{H}_3\text{O}^+]$ sont liées \emph{via} la  constante d'autoprotolyse de l'eau : \vspace{-1mm}
\[\vspace{-2mm}
K_e=\frac{[\text{HO}^-]\times[\text{H}_3\text{O}^+]}{(C^\circ)^2}\qquad\text{donc}\qquad\frac{[\text{H}_3\text{O}^+]}{C^\circ}=\frac{K_e\, C^\circ}{[\text{HO}^-]}.
\]
On a donc \minusSmallskip
\[
\text{pH}=-\log_{10}\left(\frac{[\text{H}_3\text{O}^+]}{C^\circ}\right)=
-\log_{10}\left(\frac{K_e C^\circ}{[\text{HO}^-]}\right)=
-\log_{10}\left(\frac{\num{1,0e-14}\times \SI{1,0}{\mol\per\liter}}{\SI{1,0e-2}{\mol\per\liter}}\right)=
\num{12}.
\]
\end{corrige}

% ************************************** %


% =============================================================================== %
\finEntrainement
% =============================================================================== %





% ********************************************************************* %
%                           ENTRAÎNEMENT              14                  % 
% ********************************************************************* %
% ================ Métadonnées sur l'entraînement ===================== %

\titreEntrainementFacultatif	{Quelques combats de concentration}
\hauteurLargeurCadreReponse		{6mm}{12mm}
\dureeResolutionFacultative		{3} % 1, 2, 3 ou 4
\typeEntrainement				{QCM} % Newton ou QCM ou AN ou vide (exercice "normal")
\basiqueEtTransversal			{Y}
\nombreColonnesQuestions		{1} % vide, 1, 2, 3, etc.
\avecPlusieursQuestions			{Y} % Y ou N
\initialisationEntrainement
% ===================================================================== %

Pour chacun des cas suivants, déterminer quelle solution possède la plus grande concentration en ions oxonium.

% =============================================================================== %
\debutEntrainement
% =============================================================================== %


% ^^^^^^^^^^^^^^^^^^^^^^^^^^^^^^^^^^^^^^ %
%               Question    a             %
% ^^^^^^^^^^^^^^^^^^^^^^^^^^^^^^^^^^^^^^ %

\begin{enonce}
Premier cas
\begin{listeQCM2Colonnes}
	\item Une solution de $\text{pH}=\num{1,0}$.
	\item Une solution de $\text{pH}=\num{2,0}$.
\end{listeQCM2Colonnes}
\end{enonce}
\reponse{\reponseA}
\begin{corrige}
Une solution à $\text{pH}=\num{1,0}$ possède une concentration en ions oxonium [H$_3$O$^+$] = $10^{-1,0}\si{\mole\per\liter}$, et une solution à $\text{pH}=\num{2,0}$ possède une concentration en ions oxonium [H$_3$O$^+$] = $10^{-2,0}\si{\mole\per\liter}$.
\end{corrige}

% ************************************** %


% ^^^^^^^^^^^^^^^^^^^^^^^^^^^^^^^^^^^^^^ %
%               Question    b             %
% ^^^^^^^^^^^^^^^^^^^^^^^^^^^^^^^^^^^^^^ %

\begin{enonce}
Deuxième cas
\begin{listeQCM2Colonnes}
	\item Une solution avec $[\mathrm{H_3O^+}]=\SI{5,0e-2}{\mole\per\liter}$.
	\item Une solution de $\text{pH}=\num{3,0}$.
\end{listeQCM2Colonnes}
\end{enonce}

\reponse{\reponseA}

\begin{corrige}
Une solution à $\text{pH}=\num{3,0}$ possède une concentration en ions oxonium $[\mathrm{H_3O^+}] = \SI{1.0e-3,0}{\mole\per\liter}$.
\end{corrige}

% ************************************** %


% ^^^^^^^^^^^^^^^^^^^^^^^^^^^^^^^^^^^^^^ %
%               Question       c          %
% ^^^^^^^^^^^^^^^^^^^^^^^^^^^^^^^^^^^^^^ %

\begin{enonce}
Troisième cas
\begin{listeQCM1Colonne}
	\item Une solution avec $[\mathrm{HO^-}]=\SI{2,0e-2}{\mole\per\liter}$.
	\item Une solution avec $[\mathrm{HO^-}]=\SI{8,0e-2}{\mole\per\liter}$.
\end{listeQCM1Colonne}
\medskip
\end{enonce}
\reponse{\reponseA}

\begin{corrige}
Les concentrations $[\text{HO}^-]$ et $[\text{H}_3\text{O}^+]$ sont liées \emph{via} la  constante d'autoprotolyse de l'eau : \vspace{-2mm}
\[
K_e=\frac{[\text{HO}^-]\times[\text{H}_3\text{O}^+]}{(C^\circ)^2}
\qquad\text{donc}\qquad
[\text{H}_3\text{O}^+]=\frac{K_e\cdot (C^\circ)^2}{[\text{HO}^-]}. \vspace{-3mm}
\]
On a donc :
\begin{itemize}
\item
pour la solution \reponseA :  $[\mathrm{H_3O^+}]=\SI{5,0e-13}{\mole\per\liter}$.
\item
Pour la solution \reponseB : $[\mathrm{H_3O^+}]=\SI{1.25e-13}{\mole\per\liter}$.
\end{itemize}
C'est donc la solution \reponseA{} qui est la plus concentrée en ions oxonium.
\end{corrige}

% ************************************** %


% ^^^^^^^^^^^^^^^^^^^^^^^^^^^^^^^^^^^^^^ %
%               Question     d            %
% ^^^^^^^^^^^^^^^^^^^^^^^^^^^^^^^^^^^^^^ %

\begin{enonce}
Quatrième cas
\begin{listeQCM1Colonne}
	\item Une solution avec $[\mathrm{HO^-}]=\SI{1,0e-1}{\mole\per\liter}$.
	\item Une solution de $\text{pH}=\num{9.0}$.
\end{listeQCM1Colonne}
\medskip
\end{enonce}

\reponse{\reponseB}

\begin{corrige}
Les concentrations $[\text{HO}^-]$ et $[\text{H}_3\text{O}^+]$ sont liées \emph{via} la  constante d'autoprotolyse de l'eau : \vspace{-2mm}
\[
K_e=\frac{[\text{HO}^-]\times[\text{H}_3\text{O}^+]}{(C^\circ)^2}
\qquad\text{donc}\qquad
[\text{H}_3\text{O}^+]=\frac{K_e\cdot (C^\circ)^2}{[\text{HO}^-]}.
\]
On a donc $\mathrm{[H_3O^+]}=\SI{1.0e-13}{\mole\per\liter}$ pour la solution \reponseA.

Quant à la solution de $\text{pH}=\num{9,0}$ : sa concentration en ions oxonium est $\mathrm{[H_3O^+]}=\SI{1.0e-9,0}{\mole\per\liter}$.

C'est donc la solution \reponseB{} qui est la plus concentrée en ions oxonium.
\end{corrige}

% ************************************** %

% =============================================================================== %
\finEntrainement
% =============================================================================== %


\newpage



% ********************************************************************* %
%                           ENTRAÎNEMENT              15                  % 
% ********************************************************************* %
% ================ Métadonnées sur l'entraînement ===================== %

\titreEntrainementFacultatif	{Constante d'acidité}
\hauteurLargeurCadreReponse		{10mm}{5cm}
\dureeResolutionFacultative		{2} % 1, 2, 3 ou 4
\typeEntrainement				{} % Newton ou QCM ou AN ou vide (exercice "normal")
\nombreColonnesQuestions		{1} % vide, 1, 2, 3, etc.
\avecPlusieursQuestions			{Y} % Y ou N
\initialisationEntrainement
% ===================================================================== %

On considère le couple $\mathrm{NH_4^+/NH_3}$. 

Sa constante d'acidité $K_A$  est la constante d'équilibre de la réaction : 
\[
\mathrm{NH_4^+}_\text{(aq)}+
\mathrm{H_2O}_\text{($\ell$)}=
\mathrm{NH_3}_\text{(aq)}
+\mathrm{H_3O^+}_\text{(aq)}.
\]
On donne $K_A=10^{\np{-9.2}}$ à \SI{25}{\celsius}.

% =============================================================================== %
\debutEntrainement
% =============================================================================== %


% ************************************** %

% ^^^^^^^^^^^^^^^^^^^^^^^^^^^^^^^^^^^^^^ %
%               Question     a          %
% ^^^^^^^^^^^^^^^^^^^^^^^^^^^^^^^^^^^^^^ %

\begin{enonce}
À l'aide de la loi d'action de masse, exprimer le pH en fonction de $\mathrm{p}K_A=-\log_{10}(K_A)$  ainsi que des concentrations $\mathrm{[NH_4^+]}$ et $\mathrm{[NH_3]}$. 
\minusMedskip

\end{enonce}

\reponse{$\text{pH}=\text{p}K_A+\log_{10}\left(\frac{[\text{NH}_3]}{ [\text{NH}^+_4]}\right)$}

\begin{corrigeNewpage}
La concentration en ions oxonium en solution est liée à la concentration en NH$^+_4$ et NH$_3$ \emph{via} la constante d'acidité du couple (NH$^+_4$/NH$_3$) : 
\[
K_A=\frac{[\text{NH}_3]\times[\text{H}_3\text{O}^+]}{[\text{NH}^+_4]\times C^\circ}\quad\text{d'où}\quad
\frac{[\text{H}_3\text{O}^+]}{C^\circ}=\frac{K_A [\text{NH}^+_4]}{[\text{NH}_3]}.
\]
On retrouve ainsi la formule d'Henderson : 
\[
\text{pH}=-\log_{10}\left(\frac{[\text{H}_3\text{O}^+]}{C^\circ}\right)=-\log_{10}\left(\frac{K_A[\text{NH}^+_4]}{[\text{NH}_3]}\right)=
\text{p}K_A+\log_{10}\left(\frac{[\text{NH}_3]}{ [\text{NH}^+_4]}\right).
\]
\end{corrigeNewpage}

% ************************************** %



\begin{enonce}
Sachant qu'on a $\mathrm{[NH_4^+]}=\SI{2,0e-3}{\mole\per\liter}$ et $\mathrm{[NH_3]}=\SI{1,0e-3}{\mole\per\liter}$, calculer le pH de la solution.

\end{enonce}
\reponse{$\num{8,9}$}
\begin{corrige}
Numériquement, on trouve 
\[
\text{pH}=\text{p}K_A+\log_{10}\left(\frac{[\text{NH}_3]}{ [\text{NH}^+_4]}\right)=
\num{9.2}+\log_{10}\left(\frac{\SI{1,0e-3}{\mole\per\liter}}{\SI{2,0e-3}{\mole\per\liter}}\right)=
\num{8,9}.
\]
\end{corrige}




% =============================================================================== %
\finEntrainement
% =============================================================================== %




% ********************************************************************* %
%                           ENTRAÎNEMENT        16                        % 
% ********************************************************************* %
% ================ Métadonnées sur l'entraînement ===================== %

\titreEntrainementFacultatif	{Équilibre acido-basique}
\hauteurLargeurCadreReponse		{12mm}{8cm}
\dureeResolutionFacultative		{3} % 1, 2, 3 ou 4
\typeEntrainement				{} % Newton ou QCM ou AN ou vide (exercice "normal")
\nombreColonnesQuestions		{1} % vide, 1, 2, 3, etc.
\avecPlusieursQuestions			{Y} % Y ou N
\initialisationEntrainement
% ===================================================================== %

% =============================================================================== %
\debutEntrainement
% =============================================================================== %

On introduit un volume $V=\SI{20,0}{m\liter}$ d'une solution d'acide éthanoïque $\text{CH}_3\text{COOH}$ à la concentration $C= \SI{2,00e-3}{\mol\per\liter}$ dans un bécher contenant un volume $V'=\SI{20,0}{m\liter}$ d'eau distillée. 

Un équilibre s'établit selon l'équation de réaction :
\[
\text{CH}_3\text{COOH}_\text{(aq)}+\text{H}_2\text{O}_\text{($\ell$)}=\text{CH}_3\text{COO}^-_\text{(aq)}+\text{H}_3\text{O}^+_\text{(aq)}
\]
La constante d'équilibre de cette réaction est $K_A=10^{-4,8}$ à la température de l'expérience.

% ^^^^^^^^^^^^^^^^^^^^^^^^^^^^^^^^^^^^^^ %
%               Question    a             %
% ^^^^^^^^^^^^^^^^^^^^^^^^^^^^^^^^^^^^^^ %

\begin{enonce}
Établir l'équation du second degré vérifiée par l'avancement volumique $\xi_v$ à l'état final d'équilibre

\end{enonce}
\reponse{$\xi^2_v+K_AC^\circ\, \xi_v-K_AC_1C^\circ=0$}

\begin{corrige}
On écrit un tableau d'avancement pour cette réaction, où $\xi_v$ représente l'avancement volumique : 
\begin{center}%Essai avec \frac pour voir si c'ets plus lisible
\renewcommand{\arraystretch}{1.2}
\begin{tabular}{|l|c|c||c|c|}
  \hline
&
%
\multicolumn{1}{c@{\protect\makebox[0cm]{$\mathrm{+}$}}}{%
~~~~$\mathrm{CH_3COOH_\text{(aq)}}$~~~~}                   &
\multicolumn{1}{c@{\protect\makebox[0cm]{$=$}}}{%
~~~~$\mathrm{H_2O_\text{($\ell$)}}$~~~~}                        &
\multicolumn{1}{c@{\protect\makebox[0cm]{$\mathrm{+}$}}}{%
~~~~$\mathrm{CH_3COO^-_\text{(aq)}}$~~~~}                        &
\multicolumn{1}{c|}{%
~~~~$\mathrm{H_3O^+_\text{(aq)}}$~~~}              \\ \hline
%
\textbf{\'Etat initial}        & $C_1$                       &
\text{excès}                 & 0                       &
0                                             \\ \hline
%
\textbf{\'Etat final} & $C_1-\xi_v$                       &
\text{excès}                 & $\xi_v$                       &
$\xi_v$                                            \\ \hline
\end{tabular}
\end{center}
avec $C_1=\dfrac{CV}{V+V'}=\SI{1,00e-3}{\mol\per\liter}$.

À l'équilibre, d'après la loi d'action de masse, on a :
\[
Q_\text{eq}=K_A=\frac{a(\text{CH}_3\text{COO}^-)_\text{eq}\times a(\text{H}_3\text{O}^+)_\text{eq}}{a(\text{CH}_3\text{COOH})_\text{eq}\times a(\text{H}_2\text{O})_\text{eq}}.
\]
En solution aqueuse diluée, on remplace les activités par leurs expressions. On obtient
\[
K_A=\frac{[\text{CH}_3\text{COO}^-]_\text{eq}\times [\text{H}_3\text{O}^+]_\text{eq}}{[\text{CH}_3\text{COOH}]_\text{eq}\,C^\circ}.
\]
Ensuite, on remplace les concentrations par leurs expressions trouvées dans le tableau d'avancement. Il vient : 
\[
K_A=\frac{\xi^2_v}{(C_1-\xi_v)C^\circ}
\qquad\text{donc}\qquad
\xi^2_v+K_AC^\circ\, \xi_v-K_AC_1C^\circ=0.
\]
\end{corrige}

% ************************************** %


\hauteurLargeurCadreReponse		{9mm}{5cm}


% ^^^^^^^^^^^^^^^^^^^^^^^^^^^^^^^^^^^^^^ %
%               Question    b             %
% ^^^^^^^^^^^^^^^^^^^^^^^^^^^^^^^^^^^^^^ %

\begin{enonce}
Calculer $\mathrm{[CH_3COOH]}_\text{eq}$ à l'équilibre.
\end{enonce}

\reponse{\SI{8.8e-4}{\mole\per\liter}}

\begin{corrige}
La résolution du polynôme du second degré obtenu à la question précédente donne deux solutions : 
\[
\xi_{v,1}= \SI{1,2e-4}{\mole}
\quad\text{et}\quad
\xi_{v,2}=\SI{-1,3e-4}{\mole}.
\]
Le quotient de réaction à l'instant initial vaut $Q_i=0$ (il n'y a pas de produits à l'instant initial). 

Ainsi, on a $Q_i<K_A$ : la réaction se produit dans le sens direct. L'avancement doit donc être positif et on a $\xi_v=\xi_{v,1}$.

Ainsi, à l'équilibre $[\text{CH}_3\text{COOH}]=C_1-\xi_{v}=\SI{8.8e-4}{\mole\per\liter}$.
\end{corrige}

% ************************************** %


% ^^^^^^^^^^^^^^^^^^^^^^^^^^^^^^^^^^^^^^ %
%               Question    c             %
% ^^^^^^^^^^^^^^^^^^^^^^^^^^^^^^^^^^^^^^ %

\begin{enonce}
En déduire le pH de la solution à l'équilibre.
\end{enonce}

\reponse{\num{3.9}}

\begin{corrige}
La concentration en ions $\text{H}_3\text{O}^+$ est égale à $\xi_v$, ainsi le pH se calcule par $\text{pH}=-\log_{10}\left(\dfrac{\xi_v}{C^\circ}\right)=\num{3.9}$.
\end{corrige}

% ************************************** %

% =============================================================================== %
\finEntrainement
% =============================================================================== %






% ---------------------------------------------------------------------------- %
%                       Affichage des réponses mélangées                       %
% ---------------------------------------------------------------------------- %
\afficheReponsesMelangees
% ---------------------------------------------------------------------------- %

%%%%%%%%%%%%%%%%%%%%%%%%%%%%%%%%%%%%%%%%%%%%%%%%%%%%%%%%%%%%%%%%%%%%%%%%%%%%%%%%%%%%%%%%%%%%%%
\finFicheEntrainement                                                                        %
%%%%%%%%%%%%%%%%%%%%%%%%%%%%%%%%%%%%%%%%%%%%%%%%%%%%%%%%%%%%%%%%%%%%%%%%%%%%%%%%%%%%%%%%%%%%%%


% ======================================================= % 
% ============== Gestion de la compilation ============== %
% ======================================================= % 
\ifdefined\mainIsLoaded\else
\printReponsesEtCorriges
\end{document}\fi
% ======================================================= % 
% ======================================================= % 
